\documentclass[border=10pt]{standalone}
\usepackage{luatexja-fontspec}
\usepackage{tikz}
\usetikzlibrary{positioning, arrows.meta, calc, patterns.meta, shapes.geometric}

% ヒラギノ丸ゴシックを指定
\setmainjfont{HiraMaruProN-W4}

\begin{document}
\begin{tikzpicture}[
    thick,                  % 全体の線を太めにする
    >=Stealth,              % 綺麗な矢じりを使用
    font=\large,            % フォントサイズを一括で大きくする
    window_box/.style={draw=cyan!80!black, line width=1.2pt, rectangle, minimum height=2.5cm, anchor=south west},
    window_curve/.style={draw=green!60!black, line width=1.2pt, smooth},
    text_green/.style={text=green!60!black, font=\small, align=center},
    hatch_cyan/.style={pattern={Lines[angle=45, distance=6pt, line width=0.5pt]}, pattern color=cyan!30!white},
    hatch_green/.style={pattern={Lines[angle=-45, distance=4pt, line width=0.5pt]}, pattern color=green!30!white}
]

    % =========================================================================
    % --- 上段：全体の信号波形と窓関数による切り出し（再構成） ---
    % =========================================================================
    
    % ベースとなる時間軸
    \draw [->] (-3, 1.25) -- (11.5, 1.25) node[right] {$t$};

    % 全体の長い信号（グレーの背景線）
    \draw [black!40, line width=1.0pt, smooth, tension=0.5] plot coordinates {
        (-2.5, 1.2) (-2.0, 0.7) (-1.5, 1.6) (-1.0, 0.9) (-0.5, 1.9)
        (0.0, 1.2) (0.5, 2.0) (1.0, 0.5) (1.5, 1.5) (2.0, 0.8)
        (2.5, 1.4) (3.0, 1.0) (3.5, 2.3) (4.0, 0.4) (4.5, 1.8)
        (5.0, 0.7) (5.5, 1.3) (6.0, 1.0) (6.5, 1.6) (7.0, 0.9)
        (7.5, 1.9) (8.0, 0.6) (8.5, 2.4) (9.0, 0.5) (9.5, 1.5)
        (10.0, 1.1) (10.5, 1.3) (11.0, 0.8)
    };

    % -------------------------------------------------------------------------
    % --- 窓関数1（左側の区間：ここだけハッチングを残します） ---
    % -------------------------------------------------------------------------
    % 1. 水色枠の背景ハッチング
    \path [hatch_cyan] (0, 0) rectangle (4.0, 2.5);

    % 2. 窓関数（緑の太い実線）の内部ハッチング
    \path [hatch_green] plot [domain=0:4, samples=50] (\x, {1.25 + 1.0 * sin(\x * 180 / 4)})
        -- plot [domain=4:0, samples=50] (\x, {1.25 - 1.0 * sin(\x * 180 / 4)}) -- cycle;

    % 3. 枠線とレンズ曲線の清書
    \node [window_box, minimum width=4.0cm] (wbox1) at (0, 0) {};
    \draw [window_curve] plot [domain=0:4, samples=50] (\x, {1.25 + 1.0 * sin(\x * 180 / 4)});
    \draw [window_curve] plot [domain=0:4, samples=50] (\x, {1.25 - 1.0 * sin(\x * 180 / 4)});


    % -------------------------------------------------------------------------
    % --- 窓関数2（右側の区間：ハッチングを削除し、枠線と曲線のみ描画） ---
    % -------------------------------------------------------------------------
    \node [window_box, minimum width=4.0cm] (wbox2) at (2.5, 0) {};
    \draw [window_curve] plot [domain=2.5:6.5, samples=50] (\x, {1.25 + 1.0 * sin((\x-2.5) * 180 / 4)});
    \draw [window_curve] plot [domain=2.5:6.5, samples=50] (\x, {1.25 - 1.0 * sin((\x-2.5) * 180 / 4)});


    % 区間の重ね合わせ（ずらし）を示す中央の破線
    \draw [cyan!80!black, dashed, line width=1.0pt] (2.5, 0) -- (2.5, 2.5);
    \draw [cyan!80!black, dashed, line width=1.0pt] (4.0, 0) -- (4.0, 2.5);


    % =========================================================================
    % --- 下段：切り出された断片の拡大表示（背景ハッチングを削除） ---
    % =========================================================================
    
    % 拡大表示用の軸
    \draw [->] (-1.0, -3.0) -- (6.0, -3.0) node[right] {$t$};
    \draw (-1.0, -3.0) node[below] {$0$};

    % 切り出され、端が減衰した断片信号
    \draw [black, line width=2.0pt, smooth, tension=0.5] plot coordinates {
        (0.0, -3.0) (0.4, -2.6) (0.8, -3.3) (1.2, -2.4) (1.6, -3.2)
        (2.1, -1.8) (2.5, -1.2) (2.9, -2.0) (3.5, -3.2) (4.2, -2.5)
        (4.6, -2.6) (5.0, -3.0)
    };

    % --- 上段から下段へのズーム引き出し線 ---
    \draw [cyan!70!black, line width=1.0pt, dashed] (wbox1.south west) -- (-1.0, -1.75);
    \draw [cyan!70!black, line width=1.0pt, dashed] (wbox1.south east) -- (5.0, -1.75);

    % --- 不連続性を軽減する端部の丸囲み注釈 ---
    \node [draw=green!60!black, circle, minimum size=1.2cm, inner sep=0pt, line width=1.2pt, fill=white, fill opacity=0.8] (circle_L) at (0.0, -3.0) {};
    \node [text_green, below left=0.1cm and 0.1cm of circle_L] {不連続性を\\軽減};

    \node [draw=green!60!black, circle, minimum size=1.2cm, inner sep=0pt, line width=1.2pt, fill=white, fill opacity=0.8] (circle_R) at (5.0, -3.0) {};


    % =========================================================================
    % --- 右側：箇条書きの説明文 ---
    % =========================================================================
    \matrix [draw=none, column sep=0.5cm, anchor=north west] at (6.5, -1.2) {
        \node [align=left, font=\normalsize, text width=6.5cm, yshift=0.3cm] {
            ● 各区分の中で\\
            \hspace{1em} デジタル信号処理 (フィルタ処理) \\[0.8\baselineskip]
            ● \underline{処理する区分}は\\
            \hspace{1em} 少しずつずらすことが多い \\[0.8\baselineskip]
            ● 処理後に重複も含めて\\
            \hspace{1em} 長い系列に（再構成して）戻す
        }; \\
    };

\end{tikzpicture}
\end{document}

\documentclass[border=10pt]{standalone}
\usepackage{luatexja-fontspec}
\usepackage{tikz}
\usetikzlibrary{positioning, arrows.meta, calc, patterns.meta, shapes.geometric}

% ヒラギノ丸ゴシックを指定
\setmainjfont{HiraMaruProN-W4}

\begin{document}
\begin{tikzpicture}[
    thick,                  % 全体の線を太めにする
    >=Stealth,              % 綺麗な矢じりを使用
    font=\large,            % フォントサイズを一括で大きくする
    window_box/.style={draw=cyan!80!black, line width=1.2pt, rectangle, minimum height=2.5cm, anchor=south west},
    window_curve/.style={draw=green!60!black, line width=1.2pt, smooth},
    text_green/.style={text=green!60!black, font=\small, align=center},
    hatch_cyan/.style={pattern={Lines[angle=45, distance=6pt, line width=0.5pt]}, pattern color=cyan!30!white},
    hatch_green/.style={pattern={Lines[angle=-45, distance=4pt, line width=0.5pt]}, pattern color=green!30!white}
]

    % =========================================================================
    % --- 上段：全体の信号波形と窓関数による切り出し（再構成） ---
    % =========================================================================
    
    % ベースとなる時間軸
    \draw [->] (-3, 1.25) -- (11.5, 1.25) node[right] {$t$};

    % 全体の長い信号（グレーの背景線）
    \draw [black!40, line width=1.0pt, smooth, tension=0.5] plot coordinates {
        (-2.5, 1.2) (-2.0, 0.7) (-1.5, 1.6) (-1.0, 0.9) (-0.5, 1.9)
        (0.0, 1.2) (0.5, 2.0) (1.0, 0.5) (1.5, 1.5) (2.0, 0.8)
        (2.5, 1.4) (3.0, 1.0) (3.5, 2.3) (4.0, 0.4) (4.5, 1.8)
        (5.0, 0.7) (5.5, 1.3) (6.0, 1.0) (6.5, 1.6) (7.0, 0.9)
        (7.5, 1.9) (8.0, 0.6) (8.5, 2.4) (9.0, 0.5) (9.5, 1.5)
        (10.0, 1.1) (10.5, 1.3) (11.0, 0.8)
    };

    % -------------------------------------------------------------------------
    % --- 窓関数1（左側の区間：ここを下段へクローズアップ） ---
    % -------------------------------------------------------------------------
    % 1. 水色枠の背景ハッチング
    \path [hatch_cyan] (0, 0) rectangle (4.0, 2.5);

    % 2. 窓関数（緑の太い実線）の内部ハッチング
    \path [hatch_green] plot [domain=0:4, samples=50] (\x, {1.25 + 1.0 * sin(\x * 180 / 4)})
        -- plot [domain=4:0, samples=50] (\x, {1.25 - 1.0 * sin(\x * 180 / 4)}) -- cycle;

    % 3. 枠線とレンズ曲線の清書
    \node [window_box, minimum width=4.0cm] (wbox1) at (0, 0) {};
    \draw [window_curve] plot [domain=0:4, samples=50] (\x, {1.25 + 1.0 * sin(\x * 180 / 4)});
    \draw [window_curve] plot [domain=0:4, samples=50] (\x, {1.25 - 1.0 * sin(\x * 180 / 4)});


    % -------------------------------------------------------------------------
    % --- 窓関数2（右側の区間：少し重ねて配置） ---
    % -------------------------------------------------------------------------
    % 1. 水色枠の背景ハッチング
    \path [hatch_cyan] (2.5, 0) rectangle (6.5, 2.5);

    % 2. 窓関数内部のハッチング
    \path [hatch_green] plot [domain=2.5:6.5, samples=50] (\x, {1.25 + 1.0 * sin((\x-2.5) * 180 / 4)})
        -- plot [domain=6.5:2.5, samples=50] (\x, {1.25 - 1.0 * sin((\x-2.5) * 180 / 4)}) -- cycle;

    % 3. 枠線とレンズ曲線の清書
    \node [window_box, minimum width=4.0cm] (wbox2) at (2.5, 0) {};
    \draw [window_curve] plot [domain=2.5:6.5, samples=50] (\x, {1.25 + 1.0 * sin((\x-2.5) * 180 / 4)});
    \draw [window_curve] plot [domain=2.5:6.5, samples=50] (\x, {1.25 - 1.0 * sin((\x-2.5) * 180 / 4)});


    % 区間の重ね合わせ（ずらし）を示す中央の破線
    \draw [cyan!80!black, dashed, line width=1.0pt] (2.5, 0) -- (2.5, 2.5);
    \draw [cyan!80!black, dashed, line width=1.0pt] (4.0, 0) -- (4.0, 2.5);


    % =========================================================================
    % --- 下段：切り出された断片の拡大表示 ---
    % =========================================================================
    
    % 拡大表示用のズーム引き出し領域背景
    \path [hatch_cyan] (-1.0, -4.25) rectangle (5.0, -1.75);

    % 拡大表示用の軸
    \draw [->] (-1.0, -3.0) -- (6.0, -3.0) node[right] {$t$};
    \draw (-1.0, -3.0) node[below] {$0$};

    % 切り出され、端が減衰した断片信号
    \draw [black, line width=2.0pt, smooth, tension=0.5] plot coordinates {
        (0.0, -3.0) (0.4, -2.6) (0.8, -3.3) (1.2, -2.4) (1.6, -3.2)
        (2.1, -1.8) (2.5, -1.2) (2.9, -2.0) (3.5, -3.2) (4.2, -2.5)
        (4.6, -2.6) (5.0, -3.0)
    };

    % --- 上段から下段へのズーム引き出し線 ---
    \draw [cyan!70!black, line width=1.0pt, dashed] (wbox1.south west) -- (-1.0, -1.75);
    \draw [cyan!70!black, line width=1.0pt, dashed] (wbox1.south east) -- (5.0, -1.75);

    % --- 不連続性を軽減する端部の丸囲み注釈 ---
    \node [draw=green!60!black, circle, minimum size=1.2cm, inner sep=0pt, line width=1.2pt, fill=white, fill opacity=0.8] (circle_L) at (0.0, -3.0) {};
    \node [text_green, below left=0.1cm and 0.1cm of circle_L] {不連続性を\\軽減};

    \node [draw=green!60!black, circle, minimum size=1.2cm, inner sep=0pt, line width=1.2pt, fill=white, fill opacity=0.8] (circle_R) at (5.0, -3.0) {};


    % =========================================================================
    % --- 右側：箇条書きの説明文 ---
    % =========================================================================
    \matrix [draw=none, column sep=0.5cm, anchor=north west] at (6.5, -1.2) {
        \node [align=left, font=\normalsize, text width=6.5cm, yshift=0.3cm] {
            ● 各区分の中で\\
            \hspace{1em} デジタル信号処理 (フィルタ処理) \\[0.8\baselineskip]
            ● \underline{処理する区分}は\\
            \hspace{1em} 少しずつずらすことが多い \\[0.8\baselineskip]
            ● 処理後に重複も含めて\\
            \hspace{1em} 長い系列に（再構成して）戻す
        }; \\
    };

\end{tikzpicture}
\end{document}

\documentclass[border=10pt]{standalone}
\usepackage{luatexja-fontspec}
\usepackage{tikz}
\usetikzlibrary{positioning, arrows.meta, calc, patterns.meta, shapes.geometric}

% ヒラギノ丸ゴシックを指定
\setmainjfont{HiraMaruProN-W4}

\begin{document}
\begin{tikzpicture}[
    thick,                  % 全体の線を太めにする
    >=Stealth,              % 綺麗な矢じりを使用
    font=\large,            % フォントサイズを一括で大きくする
    % 個別スタイル定義
    window_box/.style={draw=cyan!80!black, line width=1.2pt, rectangle, minimum height=2.5cm, anchor=south west},
    window_curve/.style={draw=green!60!black, line width=1.2pt, smooth},
    text_green/.style={text=green!60!black, font=\small, align=center},
    
    % 【新設】ハッチングテクスチャのスタイル定義
    % 1. 切り出し領域用（水色の粗い斜線）
    hatch_cyan/.style={
        pattern={Lines[angle=45, distance=6pt, line width=0.5pt]},
        pattern color=cyan!30!white
    },
    % 2. 窓関数内部用（緑色の細かい逆斜線：重ね合わせで格子状に見えないよう角度を変更）
    hatch_green/.style={
        pattern={Lines[angle=-45, distance=4pt, line width=0.5pt]},
        pattern color=green!30!white
    }
]

    % =========================================================================
    % --- 上段：全体の信号波形と窓関数による切り出し（再構成） ---
    % =========================================================================
    
    % ベースとなる時間軸
    \draw [->] (-3, 1.25) -- (11.5, 1.25) node[right] {$t$};

    % 全体の長い信号（グレーの背景線）
    \draw [black!40, line width=1.0pt, smooth, tension=0.5] plot coordinates {
        (-2.5, 1.2) (-2.0, 0.7) (-1.5, 1.6) (-1.0, 0.9) (-0.5, 1.9)
        (0.0, 1.2) (0.5, 2.0) (1.0, 0.5) (1.5, 1.5) (2.0, 0.8)
        (2.5, 1.4) (3.0, 1.0) (3.5, 2.3) (4.0, 0.4) (4.5, 1.8)
        (5.0, 0.7) (5.5, 1.3) (6.0, 1.0) (6.5, 1.6) (7.0, 0.9)
        (7.5, 1.9) (8.0, 0.6) (8.5, 2.4) (9.0, 0.5) (9.5, 1.5)
        (10.0, 1.1) (10.5, 1.3) (11.0, 0.8)
    };

    % -------------------------------------------------------------------------
    % --- 窓関数1（左側の区間：ここを下段へクローズアップ） ---
    % -------------------------------------------------------------------------
    % 1. 水色枠の背景ハッチング
    \path [hatch_cyan] (0, 0) rectangle (4.0, 2.5);

    % 2. 窓関数（緑の太い実線）の内部ハッチング
    \path [hatch_green] plot [domain=0:4, samples=50] (\x, {1.25 + 1.0 * sin(\x * 180 / 4)})
        -- plot [domain=4:0, samples=50] (\x, {1.25 - 1.0 * sin(\x * 180 / 4)}) -- cycle;

    % 3. 枠線とレンズ曲線の清書（ハッチングの上から描画して線を際立たせる）
    \node [window_box, minimum width=4.0cm] (wbox1) at (0, 0) {};
    \draw [window_curve] plot [domain=0:4, samples=50] (\x, {1.25 + 1.0 * sin(\x * 180 / 4)});
    \draw [window_curve] plot [domain=0:4, samples=50] (\x, {1.25 - 1.0 * sin(\x * 180 / 4)});


    % -------------------------------------------------------------------------
    % --- 窓関数2（右側の区間：少し重ねて配置） ---
    % -------------------------------------------------------------------------
    % 1. 水色枠の背景ハッチング
    \path [hatch_cyan] (2.5, 0) rectangle (6.5, 2.5);

    % 2. 窓関数内部のハッチング
    \path [hatch_green] plot [domain=2.5:6.5, samples=50] (\x, {1.25 + 1.0 * sin((\x-2.5) * 180 / 4)})
        -- plot [domain=6.5:2.5, samples=50] (\x, {1.25 - 1.0 * sin((\x-2.5) * 180 / 4)}) -- cycle;

    % 3. 枠線とレンズ曲線の清書
    \node [window_box, minimum width=4.0cm] (wbox2) at (2.5, 0) {};
    \draw [window_curve] plot [domain=2.5:6.5, samples=50] (\x, {1.25 + 1.0 * sin((\x-2.5) * 180 / 4)});
    \draw [window_curve] plot [domain=2.5:6.5, samples=50] (\x, {1.25 - 1.0 * sin((\x-2.5) * 180 / 4)});


    % 区間の重ね合わせ（ずらし）を示す中央の破線
    \draw [cyan!80!black, dashed, line width=1.0pt] (2.5, 0) -- (2.5, 2.5);
    \draw [cyan!80!black, dashed, line width=1.0pt] (4.0, 0) -- (4.0, 2.5);


    % =========================================================================
    % --- 下段：切り出された断片の拡大表示 ---
    % =========================================================================
    
    % 拡大表示用のズーム引き出し領域背景（下段の軸の範囲に対応するハッチング）
    % これにより、上の「水色枠1」の中身がそのまま下にスライドしてきたことが強調されます
    \path [hatch_cyan] (-1.0, -4.25) rectangle (5.0, -1.75);

    % 拡大表示用の軸
    \draw [->] (-1.0, -3.0) -- (6.0, -3.0) node[right] {$t$};
    \draw (-1.0, -3.0) node[below] {$0$};

    % 切り出され、端が減衰した断片信号（主役なので黒の極太線に変更）
    \draw [black, line width=2.0pt, smooth, tension=0.5] plot coordinates {
        (0.0, -3.0) (0.4, -2.6) (0.8, -3.3) (1.2, -2.4) (1.6, -3.2)
        (2.1, -1.8) (2.5, -1.2) (2.9, -2.0) (3.5, -3.2) (4.2, -2.5)
        (4.6, -2.6) (5.0, -3.0)
    };

    % --- 上段から下段へのズーム引き出し線 ---
    \draw [cyan!70!black, line width=1.0pt, dashed] (wbox1.south west) -- (-1.0, -1.75);
    \draw [cyan!70!black, line width=1.0pt, dashed] (wbox1.south east) -- (5.0, -1.75);

    % --- 不連続性を軽減する端部の丸囲み注釈 ---
    \node [draw=green!60!black, circle, minimum size=1.2cm, inner sep=0pt, line width=1.2pt, fill=white, fill opacity=0.8] (circle_L) at (0.0, -3.0) {};
    \node [text_green, below left=0.1cm and 0.1cm of circle_L] {不連続性を\\軽減};

    \node [draw=green!60!black, circle, minimum size=1.2cm, inner sep=0pt, line width=1.2pt, fill=white, fill opacity=0.8] (circle_R) at (5.0, -3.0) {};


    % =========================================================================
    % --- 右側：箇条書きの説明文 ---
    % =========================================================================
    \matrix [draw=none, column sep=0.5cm, anchor=north west] at (6.5, -1.2) {
        \node [align=left, font=\normalsize, text width=6.5cm, yshift=0.3cm] {
            ● 各区分の中で\\
            \hspace{1em} デジタル信号処理 (フィルタ処理) \\[0.8\baselineskip]
            ● \underline{処理する区分}は\\
            \hspace{1em} 少しずつずらすことが多い \\[0.8\baselineskip]
            ● 処理後に重複も含めて\\
            \hspace{1em} 長い系列に（再構成して）戻す
        }; \\
    };

\end{tikzpicture}
\end{document}

\documentclass[border=10pt]{standalone}
\usepackage{luatexja-fontspec}
\usepackage{tikz}
\usetikzlibrary{positioning, arrows.meta, calc, shapes.geometric}

% ヒラギノ丸ゴシックを指定
\setmainjfont{HiraMaruProN-W4}

\begin{document}
\begin{tikzpicture}[
    thick,                  % 全体の線を太めにする
    >=Stealth,              % 綺麗な矢じりを使用
    font=\large,            % フォントサイズを一括で大きくする
    % 個別スタイル定義
    window_box/.style={draw=cyan!80!black, line width=1.2pt, rectangle, minimum height=2.5cm, anchor=south west},
    window_curve/.style={draw=green!60!black, line width=1.0pt, smooth},
    text_green/.style={text=green!60!black, font=\small, align=center}
]

    % =========================================================================
    % --- 上段：全体の信号波形と窓関数による切り出し（再構成） ---
    % =========================================================================
    
    % ベースとなる時間軸
    \draw [->] (-3, 1.25) -- (11.5, 1.25) node[right] {$t$};

    % 全体の長い信号（適度な凸凹を持つ長い波形）
    \draw [black!70, line width=1.0pt, smooth, tension=0.5] plot coordinates {
        (-2.5, 1.2) (-2.0, 0.7) (-1.5, 1.6) (-1.0, 0.9) (-0.5, 1.9)
        (0.0, 1.2) (0.5, 2.0) (1.0, 0.5) (1.5, 1.5) (2.0, 0.8)
        (2.5, 1.4) (3.0, 1.0) (3.5, 2.3) (4.0, 0.4) (4.5, 1.8)
        (5.0, 0.7) (5.5, 1.3) (6.0, 1.0) (6.5, 1.6) (7.0, 0.9)
        (7.5, 1.9) (8.0, 0.6) (8.5, 2.4) (9.0, 0.5) (9.5, 1.5)
        (10.0, 1.1) (10.5, 1.3) (11.0, 0.8)
    };

    % --- 窓関数1（左側の区間） ---
    % 水色の枠（切り出し区間）
    \node [window_box, minimum width=4.0cm] (wbox1) at (0, 0) {};
    % 緑の窓関数（ハニング窓・ハミング窓を模したなだらかな山型）
    % 上の窓関数
    \draw [window_curve] plot [domain=0:4, samples=50] (\x, {1.25 + 1.0 * sin(\x * 180 / 4)});
    % 下の窓関数（上下対称のガイド）
    \draw [window_curve] plot [domain=0:4, samples=50] (\x, {1.25 - 1.0 * sin(\x * 180 / 4)});

    % --- 窓関数2（右側の区間：少し重ねて配置） ---
    % 水色の枠（x=2.5から少しオーバーラップさせて配置）
    \node [window_box, minimum width=4.0cm] (wbox2) at (2.5, 0) {};
    % 上の窓関数（平行移動）
    \draw [window_curve] plot [domain=2.5:6.5, samples=50] (\x, {1.25 + 1.0 * sin((\x-2.5) * 180 / 4)});
    % 下の窓関数
    \draw [window_curve] plot [domain=2.5:6.5, samples=50] (\x, {1.25 - 1.0 * sin((\x-2.5) * 180 / 4)});

    % 区間の重ね合わせ（ずらし）を示す中央の破線
    \draw [cyan!80!black, dashed] (2.5, 0) -- (2.5, 2.5);
    \draw [cyan!80!black, dashed] (4.0, 0) -- (4.0, 2.5);


    % =========================================================================
    % --- 下段：切り出された断片の拡大表示 ---
    % =========================================================================
    
    % 拡大表示用の軸
    \draw [->] (-1.0, -3.0) -- (6.0, -3.0) node[right] {$t$};
    \draw (-1.0, -3.0) node[below] {$0$};

    % 切り出され、端が減衰した断片信号（不連続性が緩和された波形）
    % 両端（x=0 と x=5.0）でなめらかに0（中心軸）に収束する波形
    \draw [black, line width=1.5pt, smooth, tension=0.5] plot coordinates {
        (0.0, -3.0) (0.4, -2.6) (0.8, -3.3) (1.2, -2.4) (1.6, -3.2)
        (2.1, -1.8) (2.5, -1.2) (2.9, -2.0) (3.5, -3.2) (4.2, -2.5)
        (4.6, -2.6) (5.0, -3.0)
    };

    % --- 上段から下段へのズーム引き出し線 ---
    \draw [cyan!70!black, line width=0.8pt] (wbox1.south west) -- (-1.0, -2.0);
    \draw [cyan!70!black, line width=0.8pt] (wbox1.south east) -- (5.0, -2.0);

    % --- 不連続性を軽減する端部の丸囲み注釈 ---
    % 左端の緑丸と注釈
    \node [draw=green!60!black, circle, minimum size=1.2cm, inner sep=0pt] (circle_L) at (0.0, -3.0) {};
    \node [text_green, below left=0.1cm and 0.1cm of circle_L] {不連続性を\\軽減};

    % 右端の緑丸
    \node [draw=green!60!black, circle, minimum size=1.2cm, inner sep=0pt] (circle_R) at (5.0, -3.0) {};


    % =========================================================================
    % --- 右側：箇流書きの説明文 ---
    % =========================================================================
    \matrix [draw=none, column sep=0.5cm, anchor=north west] at (6.5, -1.2) {
        \node [align=left, font=\normalsize, text width=6.5cm, yshift=0.3cm] {
            ● 各区分の中で\\
            \hspace{1em} デジタル信号処理 (フィルタ処理) \\[0.8\baselineskip]
            ● \underline{処理する区分}は\\
            \hspace{1em} 少しずつずらすことが多い \\[0.8\baselineskip]
            ● 処理後に重複も含めて\\
            \hspace{1em} 長い系列に（再構成して）戻す
        }; \\
    };

\end{tikzpicture}
\end{document}

